\section*{Chapitre \ref{ch:iv1:timed-test-1} --- Épreuve chronométrée 1}
\markboth{\protect\omTocMark{}{Réponses~: Épreuve chronométrée 1}}{\protect\omTocMark{}{Réponses~: Épreuve chronométrée 1}}

\begin{solution}{iq:iv1:timed-test-1:1}
La réponse est $144/1024 = 9/64 \approx 0.141$.

Soit $a_n$ le nombre de suites de $n$ lancers sans deux piles consécutifs. Une bonne suite se termine soit par un pile, après toute bonne suite de longueur $n - 1$, soit par un pile-face, après toute bonne suite de longueur $n - 2$. D'où $a_n = a_{n-1} + a_{n-2}$\index{suite de Fibonacci}, avec $a_1 = 2$ et $a_2 = 3$. La suite vaut $2, 3, 5, 8, 13, 21, 34, 55, 89, 144$, et il y a $2^{10} = 1024$ suites équiprobables.

\iqwhy{Écrire les deux derniers lancers est la voie la plus rapide pour la plupart des dénombrements de ce type, et elle évite une longue analyse par cas.}
\end{solution}

\begin{solution}{iq:iv1:timed-test-1:2}
Le premier joueur gagne avec une probabilité de $4/9$.

Utiliser l'analyse au premier pas\index{analyse au premier pas}, avec deux états pour le joueur qui va lancer. Soit $y$ sa probabilité de gagner lorsque le lancer précédent était un pile, et $x$ lorsqu'il ne l'était pas. Si le lancer précédent était un pile, face gagne immédiatement, et pile transmet le même état à l'adversaire~:
\[
y = \tfrac12 + \tfrac12(1 - y), \quad\text{donc } y = \tfrac23.
\]
Sinon, tout lancer passe la main, pile donnant à l'adversaire le bon état~:
\[
x = \tfrac12(1 - y) + \tfrac12(1 - x), \quad\text{donc } x = \tfrac49.
\]
Le premier joueur part dans l'état $x$.

\iqwrong{Commencer n'est pas un avantage ici. Votre premier lancer ne peut pas gagner, et il peut donner à l'adversaire une chance de gagner.}
\end{solution}

\begin{solution}{iq:iv1:timed-test-1:3}
La réponse est 14 chutes, en partant de l'étage 14.

Si votre plan autorise $k$ chutes, la première chute peut se faire au plus depuis l'étage $k$~: si l'œuf se brise, le second œuf doit remonter les $k - 1$ étages inférieurs, un par un. S'il survit, il reste $k - 1$ chutes, donc le saut suivant est de $k - 1$ étages, et ainsi de suite. Ainsi $k$ chutes permettent de couvrir $k + (k - 1) + \dots + 1 = k(k + 1)/2$ étages\index{nombres triangulaires}. Comme $13 \cdot 14/2 = 91$ est insuffisant et que $14 \cdot 15/2 = 105$ suffit, $k = 14$, le premier œuf étant lâché depuis les étages 14, 27, 39, et ainsi de suite.

\iqwrong{Des pas égaux de 10, qui coûtent jusqu'à 19 chutes. Des pas égaux gaspillent le coût des premières chutes.}
\iqfollow{Et avec 3 œufs~?}{9 chutes.}
\end{solution}

\begin{solution}{iq:iv1:timed-test-1:4}
La réponse est $5/16$.

La calculer par des aires dans le carré unité\index{probabilité géométrique}. L'événement $|X - Y| > a$ est formé de deux triangles rectangles aux coins, chacun de côtés de longueur $1 - a$, d'aire totale $(1 - a)^2$. Donc $P(|X - Y| < a) = 1 - (1 - a)^2$, et
\[
P\bigl(\tfrac14 < |X - Y| < \tfrac12\bigr) = \bigl(\tfrac34\bigr)^2 - \bigl(\tfrac12\bigr)^2 = \tfrac{9}{16} - \tfrac{4}{16} = \tfrac{5}{16} ,
\]
ce qui correspond à l'aire hachurée de la \cref{fig:iv1:band}.

\begin{omfigure}
\begin{tikzpicture}[x=3.6cm,y=3.6cm]
\fill[omDef!25] (0,0.25) -- (0,0.5) -- (0.5,1) -- (0.75,1) -- cycle;
\fill[omDef!25] (0.25,0) -- (0.5,0) -- (1,0.5) -- (1,0.75) -- cycle;
\draw (0,0) rectangle (1,1);
\draw[gray,dashed] (0,0) -- (1,1);
\draw[omDef] (0,0.25) -- (0.75,1);
\draw[omDef] (0,0.5) -- (0.5,1);
\draw[omDef] (0.25,0) -- (1,0.75);
\draw[omDef] (0.5,0) -- (1,0.5);
\node[left,font=\footnotesize] at (0,0.25) {$\tfrac14$};
\node[left,font=\footnotesize] at (0,0.5) {$\tfrac12$};
\node[left,font=\footnotesize] at (0,1) {$1$};
\node[below,font=\footnotesize] at (0.25,0) {$\tfrac14$};
\node[below,font=\footnotesize] at (0.5,0) {$\tfrac12$};
\node[below,font=\footnotesize] at (1,0) {$1$};
\node[below left,font=\footnotesize] at (0,0) {$0$};
\node[right,font=\footnotesize] at (1.02,0.5) {$X$ en abscisse, $Y$ en ordonnée};
\end{tikzpicture}

\omcaption{Le carré unité avec les deux bandes hachurées où $|X - Y|$ est compris entre $\tfrac14$ et $\tfrac12$. La droite en pointillés est $X = Y$.}\label{fig:iv1:band}
\end{omfigure}
\end{solution}

\begin{solution}{iq:iv1:timed-test-1:5}
Environ $0.39$ des mois sont perdants, soit $0.386$ à trois décimales près. Toute valeur entre $0.38$ et $0.39$ convient à la main.

Sur 12 mois indépendants, la moyenne et la variance sont toutes deux multipliées par 12, donc le ratio de Sharpe annuel vaut $\sqrt{12}$ fois le ratio mensuel\index{ratio de Sharpe}. Le ratio de Sharpe mensuel vaut donc $1/\sqrt{12} \approx 0.29$. Un mois est perdant lorsqu'une normale centrée réduite passe sous $-0.29$. Au voisinage de zéro, la densité normale vaut environ $0.4$, d'où cette probabilité d'environ $0.5 - 0.29 \cdot 0.4 = 0.384$, la valeur exacte étant $0.386$.

\iqwhy{Même une stratégie de ratio de Sharpe égal à 1 perd dans près de deux mois sur cinq. Attendez-vous-y avant de juger un historique en production sur un mauvais trimestre.}
\end{solution}

\begin{solution}{iq:iv1:timed-test-1:6}
Le déterminant vaut $81/128 \approx 0.633$.

Écrire la matrice sous la forme $0.75\,I + 0.25\,\mathbf{1}\mathbf{1}^\top$, où $\mathbf{1}$ est le vecteur des uns\index{valeurs propres d'une mise à jour de rang faible}. La matrice $\mathbf{1}\mathbf{1}^\top$ a pour valeur propre 5 dans la direction de $\mathbf{1}$, et 0 sur les quatre directions orthogonales à celui-ci. Les valeurs propres sont donc $0.75 + 0.25 \cdot 5 = 2$, une fois, et $0.75$, quatre fois. Le déterminant est leur produit~: $2 \cdot (3/4)^4 = 81/128$.

\iqfollow{Jusqu'à quelle valeur peut descendre le coefficient commun hors diagonale $c$ pour que cette matrice reste une matrice de corrélation~?}{La grande valeur propre devient $1 + 4c$, qui ne doit pas être négative. Donc $c \ge -1/4$.}
\end{solution}

\begin{solution}{iq:iv1:timed-test-1:7}
La probabilité vaut $1/99$. Tous les totaux de 1 à 99 sont équiprobables.

Il s'agit d'une urne de Pólya\index{urne de Pólya}~: chaque issue ajoute un à son propre compte, et la probabilité suivante est la proportion. Démontrons l'affirmation par récurrence. Après $n$ tirs, supposons que le nombre de réussites soit uniforme sur $1, \dots, n - 1$, chaque valeur ayant donc la probabilité $1/(n - 1)$. Pour obtenir $k$ réussites sur $n + 1$ tirs, soit vous en aviez $k - 1$ et avez réussi le tir, soit vous en aviez $k$ et avez manqué le tir~:
\[
\frac{1}{n - 1}\left(\frac{k - 1}{n} + \frac{n - k}{n}\right) = \frac{1}{n - 1}\cdot\frac{n - 1}{n} = \frac1n .
\]
Donc après $n + 1$ tirs, le nombre de réussites est uniforme sur $1, \dots, n$. La propriété est vraie dès $n = 2$, et avec 100 tirs il y a 99 valeurs.

\iqwrong{Attendre un pic en 50 comme pour une binomiale. Les tirs ne sont pas indépendants, et une série initiale se prolonge.}
\end{solution}

\enlargethispage{-4\baselineskip}%
\begin{solution}{iq:iv1:timed-test-1:8}
La longueur moyenne du plus grand morceau vaut $11/18 \approx 0.611$.

Utiliser la formule des queues\index{formule de la somme des queues} avec l'inclusion--exclusion. Un morceau donné dépasse $x$ avec la probabilité $(1 - x)^2$. Deux morceaux choisis dépassent tous deux $x$ avec la probabilité $(1 - 2x)^2$ lorsque $2x < 1$, et les trois avec la probabilité $(1 - 3x)^2$ lorsque $3x < 1$. D'où
\[
P(L > x) = 3(1 - x)^2 - 3(1 - 2x)_+^2 + (1 - 3x)_+^2 .
\]
En intégrant sur $x$ de 0 à 1, on obtient $1 - \tfrac12 + \tfrac19 = \tfrac{11}{18}$.

\iqfollow{Existe-t-il une voie plus rapide~?}{Le plus petit morceau a pour moyenne $\tfrac13 \cdot \tfrac13$, le suivant est plus grand en moyenne de $\tfrac13 \cdot \tfrac12$, et le plus grand d'un supplément de $\tfrac13 \cdot 1$. Le plus grand a donc pour moyenne $\tfrac13(1 + \tfrac12 + \tfrac13) = \tfrac{11}{18}$. Avec $n$ morceaux, elle vaut $H_n/n$, où $H_n$ est la somme harmonique.}
\end{solution}

\begin{solution}{iq:iv1:timed-test-1:9}
La réponse est 7 bandelettes.

Chaque bandelette a trois issues possibles~: positive au premier tour, positive au second tour, ou jamais positive. Avec $s$ bandelettes, il y a au plus $3^s$ issues, et chaque flacon doit avoir sa propre issue\index{borne d'information}. Comme $3^6 = 729 < 1000 \le 2187 = 3^7$, six bandelettes ne suffisent pas, même avec des choix au second tour qui dépendent du premier~: si $k$ bandelettes sont positives au premier tour, il ne reste que $2^{s-k}$ issues, et $\sum_k \binom{s}{k}2^{s-k} = 3^s$ dans tous les cas.

Sept suffisent. Il faut écrire le numéro de chaque flacon en base 3 sur 7 chiffres. La bandelette $i$ teste un flacon au premier tour si le chiffre $i$ vaut 0, au second tour s'il vaut 1, et pas du tout s'il vaut 2. Les bandelettes qui réagissent, et le moment où elles le font, donnent les chiffres du flacon empoisonné.

\iqwrong{La réponse 10, issue de $2^{10} \ge 1000$, qui ignore entièrement le second tour.}
\end{solution}

\newpage
\begin{solution}{iq:iv1:timed-test-1:10}
La probabilité d'être mouillé lors d'un trajet vaut $21/370 \approx 5.68~\%$.

Il faut modéliser le nombre $i$ de parapluies à l'endroit où l'on se trouve\index{loi stationnaire}, avec $r = 3$ au total, une probabilité de pluie $p = 0.3$ et $q = 0.7$. Sans parapluie, le lieu suivant détient les $r$. Avec $i > 0$, le lieu suivant détient $r - i$ s'il ne pleut pas, et $r - i + 1$ s'il pleut et que l'on en emporte un. La résolution des équations de balance donne
\[
\pi_0 = \frac{q}{r + q}, \qquad \pi_i = \frac{1}{r + q} \quad (i = 1, \dots, r).
\]
On est mouillé lorsqu'on n'a aucun parapluie et qu'il pleut, donc la fraction vaut $p\,\pi_0 = 0.3 \cdot 0.7/3.7 = 21/370$.

\iqwrong{En prenant une loi uniforme, $\pi_0 = 1/4$, ce qui donne $0.075$. L'état zéro est plus rare que les autres, dans le rapport $q$.}
\end{solution}

\begin{solution}{iq:iv1:timed-test-1:11}
A survit avec la probabilité $25/63 \approx 0.397$, en tirant en l'air au premier tour.

Raisonnons à rebours. Dans un duel entre A et B où A tire le premier, un tour complet passe sans touche avec la probabilité $\tfrac23 \cdot \tfrac13 = \tfrac29$, donc A gagne avec la probabilité $\tfrac13/(1 - \tfrac29) = \tfrac37$.

Avec les trois encore en vie, B doit viser C, sinon C tue B immédiatement. C, s'il est appelé à tirer, vise B, le meilleur tireur. Ainsi, si A tire en l'air, B touche C avec la probabilité $2/3$, ce qui laisse à A un duel dont il tire le premier, valant $3/7$. Sinon, C tue B, et A dispose d'un tir sur C, valant $1/3$. Au total
\[
\tfrac23 \cdot \tfrac37 + \tfrac13 \cdot \tfrac13 = \tfrac{25}{63}.
\]
Viser C est moins bon~: une touche laisse à A le soin de recevoir le tir de B en premier, ce qui ne vaut que $1/7$, soit un total de $59/189$. Viser B est le pire, avec un total de $50/189$, puisque toucher B signifie que C tue A\index{induction à rebours}.

\iqwhy{Le tireur le plus faible fait le mieux en ne tirant pas, car les deux autres doivent d'abord régler leur compte entre eux.}
\end{solution}

\begin{solution}{iq:iv1:timed-test-1:12}
Il faut environ 194 lancers. Toute valeur comprise entre 190 et 200 convient.

Notons $n$ le nombre de lancers, et travaillons avec la proportion de faces\index{puissance d'un test}. Le test rejette lorsque la proportion s'éloigne de plus de $1.96 \cdot 0.5/\sqrt n$ de $0.5$. Sous l'alternative, la proportion a pour moyenne $0.6$ et pour écart type $\sqrt{0.24}/\sqrt n \approx 0.49/\sqrt n$. Pour une puissance de 80~\%, le seuil de rejet doit se situer $0.84$ écart type en dessous de $0.6$. Donc
\[
0.1 = \frac{1.96 \cdot 0.5 + 0.84 \cdot 0.49}{\sqrt n} = \frac{1.39}{\sqrt n}, \qquad n \approx 193 .
\]
En gardant davantage de décimales, on obtient 194, et la puissance binomiale exacte pour 194 lancers vaut $0.81$. Le raccourci usuel consistant à prendre $0.5$ pour les deux écarts types donne 196, ce qui convient.

\iqwrong{Utiliser $1.645$, la valeur unilatérale, donne environ 152. Ce serait la réponse à une autre question.}
\end{solution}
